\section{Introduction}
An online customer-service assistant must recover what a user means, find the relevant policy, and preserve its qualifications when answering. A follow-up such as \emph{What should I bring?} is incomplete without the previous exchange. A compound request can also involve several services, such as replacing a driver's license while updating benefit-payment details. Retrieving a plausible paragraph is therefore only one part of the task: the answer must cover each requested item without transferring a condition from one policy statement to another.

Retrieval-augmented generation (RAG; \citealt{lewis2020rag}) provides external evidence, but a fixed retrieve-and-generate pass can fail through an unresolved query, missing evidence, or an unsupported interpretation of retrieved text. These failures need different responses. Missing user information calls for clarification; missing policy evidence calls for retrieval; an incomplete answer supported by existing evidence calls for revision. Repeating all stages indiscriminately adds cost and may reproduce the same error.

We investigate a modular system with dialogue planning, configurable retrieval, evidence-conditioned generation, and an evaluator--optimizer controller. Retrieval evolved from parallel hybrid search with optional reranking to an alternative adapted Dense path without either fusion or reranking. Agents are logical roles with structured interfaces; the language-model roles share DeepSeek Flash. The controller distinguishes evidence adequacy from answer completeness and selects bounded corrective actions.

Our investigation asks three downstream questions. Does routing a compound request to separate document domains improve evidence coverage? Does generating independent domain answers help beyond pooling the same evidence? Does evaluation-driven revision increase useful answer delivery at an acceptable cost? A complementary retrieval study examines segmentation, hybrid retrieval, reranking, and encoder adaptation.

Our contributions are a runnable integration with citations and resource controls, a native retrieval study of segmentation and adaptation, a 1,000-request single-pass development evaluation, and a paired 20-request study of routing and answer composition. The small study supports domain routing more strongly than additional generation stages. Evaluator approval and citation validity remain insufficient evidence of factual correctness; we report negative findings and protocol limitations alongside gains.
